% Generated draft: requires mathematical and source review.
\documentclass[11pt]{article}
\usepackage[margin=25mm]{geometry}
\usepackage{fontspec}
\setmainfont{Latin Modern Roman}
\usepackage{amsmath,amssymb,mathtools,unicode-math}
\setmathfont{Latin Modern Math}
\usepackage{enumitem,needspace,xurl,hyperref}
\setlength{\parindent}{0pt}
\setlength{\parskip}{5pt}
\setlength{\emergencystretch}{3em}
\newcommand{\N}{\mathbb{N}}
\newcommand{\Z}{\mathbb{Z}}
\newcommand{\Q}{\mathbb{Q}}
\newcommand{\R}{\mathbb{R}}
\newcommand{\C}{\mathbb{C}}
\newcommand{\F}{\mathbb{F}}
\newcommand{\A}{\mathbb{A}}
\newcommand{\PP}{\mathbb{P}}
\newcommand{\E}{\mathbb{E}}
\newcommand{\eps}{\varepsilon}
\newcommand{\dd}{\,\mathrm{d}}
\begin{document}
\section*{UnsolvedMath ID: 2302038}
\textbf{Source problem code:} catalogue entry
\subsection{Definitions and mathematical statement}
The remaining case identified in the supplied source is whether its condition
\[
 m_0(r,f)\,M(r,f)\longrightarrow\infty\qquad(r\to\infty)
\]
can occur for a function of order one and maximal type. The record does not define $m_0(r,f)$ or state the ambient function class in this isolated excerpt. These missing definitions are essential; in particular $m_0$ is not silently replaced by the usual minimum modulus.

For contextual entire-function notation, $M(r,f)=\max_{|z|=r}|f(z)|$, and order one with maximal type means
\[
 \limsup_{r\to\infty}\frac{\log\log M(r,f)}{\log r}=1,
 \qquad
 \limsup_{r\to\infty}\frac{\log M(r,f)}{r}=\infty.
\]
Under that ambient interpretation, the question is existence of an $f$ in this class satisfying the displayed source-defined $m_0M$ behavior. The source context must confirm this interpretation and the meaning of $m_0$. The ordinary minimum modulus vanishes at the radius of any zero, so using it here requires particular care with the quoted limit and cannot be assumed merely from the notation.

The excerpt attributes to Kjellberg the implication, for $0<\alpha<1$,
\[
 \log m_0(r,f)<\cos(\phi\alpha)\log M(r,f)+O(1)
 \ \Longrightarrow\ 
 \lim_{r\to\infty}\frac{\log M(r,f)}{r^\alpha}=\beta,
 \qquad0<\beta\leq\infty.
\]
The symbol $\phi$ is not defined in the record; it is preserved literally rather than replaced by $\pi$. For $\alpha=1$, the source attributes to Hayman a corresponding conclusion with $\beta=\infty$, except for $f(z)=Ae^{Bz}$. It also reports Hayman examples satisfying the quoted $m_0M\to\infty$ condition with order $1+\epsilon$ for every $\epsilon>0$. The exact hypotheses and notation needed to reconcile these statements must be recovered before treating them as complete theorems.
\subsection{Short English statement}
Resolve the source's remaining order-one, maximal-type case for the stated minimum/maximum-type growth condition. The undefined symbol $m_0$ and cosine parameter $\phi$ prevent a standalone reconstruction; neither is silently corrected. The supplied triage says Hayman's source explicitly leaves this case open and no resolution was located. The Kjellberg and Hayman assertions are reported background, not new proofs or independently checked literature claims.
\subsection{Sources}
\begin{itemize}
\item UnsolvedMath permanent ID 2302038, source code catalogue entry, ``Research Problems in Function Theory --- Problem 2.38.'' Registry curation and metadata: Ulam AI and UnsolvedMath Contributors, CC BY 4.0; updated 2026-08-06.
\item Hayman, \emph{Research Problems in Function Theory} (2018 source update), Problem 2.38: \url{https://arxiv.org/abs/1809.07200}. Supplied source-TeX extraction accessed 2026-07-29. The metadata retains the update for status triage and marks status/release and rights review NEEDS\_REVIEW.
\item Kjellberg and Hayman are named in the source for the implications and examples, without complete bibliographic entries or theorem locators. None are inferred.
\end{itemize}
\end{document}
